\begin{tikzpicture}[
  >=Latex,
  font=\sffamily\scriptsize,
  every node/.style={align=center},
  lbl/.style={draw, thick, fill=white, inner sep=2.5pt},
  smalllbl/.style={draw, thick, fill=white, inner sep=2pt},
  wire/.style={thick},
  beam/.style={thick, dashed},
  arrow/.style={->, thick},
]

% https://www.chemguide.co.uk/analysis/masspec/masspec.GIF

  % vaporised sample inlet
  \draw[arrow] (-2.1,0.0) -- (-0.35,0.0);
  \node[anchor=east] at (-2.15,-0.0) {sample};

  % Ion source / acceleration tube
  \draw[wire] (-0.35,-0.35) rectangle (2.15,0.35);

  % Ionisation element
  \draw[wire] (0.15,-0.05) .. controls (0.35,0.25) and (0.55,0.25) .. (0.75,-0.05);
  \draw[wire] (0.75,-0.05) .. controls (0.55,-0.35) and (0.35,-0.35) .. (0.15,-0.05);

  % Acceleration slits
  \foreach \x in {1.10,1.35,1.60} {
    \draw[wire] (\x,-0.35) -- (\x,0.35);
  }

  % Labels
  \node[lbl] (ionlab) at (0.2,2.15) {IONISATION};
  \draw[arrow] (ionlab.south) -- (0.45,0.38);

  \node[lbl] (acclab) at (2.8,2.25) {ACCELERATION};
  \draw[arrow] (acclab.south west) -- (1.70,0.40);

  % Path points
  \coordinate (P0) at (2.15,0.0);
  \coordinate (P1) at (3.25,0.0);
  \coordinate (P2) at (5.10,-0.85);
  \coordinate (P3) at (6.95,-2.10);

  % Vacuum tube outline
  \draw[wire] (P0) -- (P1);
  \draw[wire] (P1) .. controls (4.10,0.00) and (4.70,-0.35) .. (P2)
              .. controls (5.70,-1.35) and (6.10,-1.70) .. (P3);

  \draw[wire] ($(P0)+(0,-0.25)$) -- ($(P1)+(0,-0.25)$);
  \draw[wire] ($(P1)+(0,-0.25)$) .. controls (4.05,-0.25) and (4.65,-0.60) .. ($(P2)+(0,-0.25)$)
              .. controls (5.60,-1.60) and (6.05,-1.95) .. ($(P3)+(0,-0.25)$);

    % ------------------------------------------------------------
    % Electromagnet poles: SAME Bezier as tube, but (i) shortened and (ii) symmetric
    
    \def\magGap{0.4}   % symmetric offset from centerline (tune 0.35..0.60)
    \def\magW{9pt}      % pole thickness (tune 8pt..12pt)
    
    % We draw only the "middle" of the bend by taking two points on the curve:
    % one near the start of deflection and one near the end.
    
    % Define two anchor points along the bend (tune these if needed)
    \coordinate (Mstart) at (3.2,-0.05);
    
    % Define matching control points for this shorter segment (tune slightly if needed)
    \coordinate (Mc1) at (3.9,-0.05);
    \coordinate (Mc2) at (4.3,-0.2);
    
    % Upper pole
    \draw[brown!55, line width=\magW, line cap=butt, line join=round]
      ($(Mstart)+(0,\magGap)$)
        .. controls ($(Mc1)+(0,\magGap)$) and ($(Mc2)+(0.15,\magGap)$) .. ($(P2)+(-0.05,\magGap)$);
    
    % Lower pole (perfect mirror: -\magGap)
    \draw[brown!55, line width=\magW, line cap=butt, line join=round]
      ($(Mstart)-(0,\magGap)$)
        .. controls ($(Mc1)-(0,\magGap)$) and ($(Mc2)-(-0.15,\magGap)$) .. ($(P2)-(0.1,\magGap)$);
    
    % Label + arrow
    \node[lbl] (magn) at (5.0,0.95) {ELECTROMAGNET};
    \draw[arrow] (magn.south) -- ($(P2)+(-0.9,0.9)$);
    \draw[arrow] (magn.south) -- ($(P2)+(-0.8,0.1)$);
    % ------------------------------------------------------------

    % Ion paths
    \draw[beam, green!60!black] (P0) -- (P1)
    .. controls (4.10,0.00) and (4.70,-0.35) .. (P2)
    .. controls (5.70,-1.35) and (6.10,-1.70) .. (P3);
    
    \draw[beam, blue] ($(P0)+(0,0.10)$) -- ($(P1)+(0,0.10)$)
    .. controls (4.05,0.10) and (4.65,-0.25) .. ($(P2)+(0,0.10)$)
    .. controls (5.60,-1.20) and (6.00,-1.55) .. ($(P3)+(0,0.10)$);
    
    \draw[beam, red] ($(P0)+(0,-0.10)$) -- ($(P1)+(0,-0.10)$)
    .. controls (4.15,-0.10) and (4.75,-0.45) .. ($(P2)+(0,-0.10)$)
    .. controls (5.80,-1.50) and (6.20,-1.85) .. ($(P3)+(0,-0.10)$);
    
    % DEFLECTION label
    \node[lbl] (deflab) at (3.15,-2) {DEFLECTION};
    \draw[arrow] (deflab.north) -- (4,-1);
    
    % Detector + electronics
    \draw[wire] ($(P3)+(0.1,-0.4)$) rectangle ($(P3)+(0.4,0.4)$);
    
    \node[lbl] (detlab) at (5.15,-3.65) {DETECTION};
    \draw[arrow] (detlab.north) -- ($(P3)+(0.1,-0.2)$);
    
    \node[smalllbl, minimum width=1.35cm, minimum height=0.55cm]
    (amp) at ($(P3)+(2.0,0.0)$) {amplifier \\ \& other electronics};
    
    \draw[arrow] ($(P3)+(0.45,-0.0)$) -- (amp.west);
\end{tikzpicture}